% gerado por tabelas_curtas.py -- não editar à mão
\begin{table}[t]
  \centering
  \caption{GuitarSet \emph{test} (58 faixas): variantes e baselines de terceiros. Acurácia em \%, média ponderada pela duração. \emph{rot.} = \texttt{root} no tempo rotulado; \emph{cob.} = cobertura; \emph{mm} e \emph{7ths} = \texttt{majmin} e \texttt{sevenths} em \texttt{chord[0]} (as demais colunas são iguais nas duas camadas).}
  \label{tab:principal}
  \footnotesize
  \begin{tabular}{lrrrrrr}
    \toprule
    Sistema & root & rot. & cob. & mm & 7ths & overs. \\
    \midrule
    \texttt{V4} & 55.9 & 55.9 & 100.0 & 48.3 & 43.6 & 75.5 \\
    \texttt{V4b} & 54.2 & 54.2 & 100.0 & 47.4 & 43.3 & 58.1 \\
    \texttt{C} & 41.5 & 65.3 & 57.1 & 39.1 & 27.3 & 59.4 \\
    \texttt{D} & 49.4 & 59.8 & 77.1 & 46.2 & 31.3 & 74.6 \\
    \texttt{E} & 47.7 & 59.1 & 75.3 & 45.7 & 31.5 & 84.2 \\
    \texttt{C0rg} & 55.0 & 55.0 & 100.0 & 49.6 & 31.9 & 55.2 \\
    \texttt{D0rg} & 57.7 & 57.7 & 100.0 & 53.5 & 34.5 & 71.3 \\
    \texttt{F0rg} & 56.9 & 56.9 & 100.0 & 52.1 & 47.1 & 73.6 \\
    \midrule
    Chordino & 66.5 & 67.4 & 98.7 & 62.5 & 41.0 & 81.8 \\
    madmom & 73.1 & 76.8 & 93.1 & 71.3 & 63.5 & 86.6 \\
    \bottomrule
  \end{tabular}
  \vspace{2pt}
  \parbox{\linewidth}{\footnotesize Fonte: \texttt{externo\_resumo\_20261117\_122651.csv}}
\end{table}
